% !TEX TS-program = pdflatex
% !TEX encoding = UTF-8 Unicode

\documentclass{beamer}

\usepackage{graphicx}

\mode<presentation>
{
  \usetheme{Warsaw}
  % or ...

  \setbeamercovered{transparent}
  % or whatever (possibly just delete it)
}


\usepackage[english]{babel}
% or whatever

\usepackage[utf8]{inputenc}
% or whatever

\usepackage{times}
\usepackage[T1]{fontenc}
\usepackage{graphicx}
% Or whatever. Note that the encoding and the font should match. If T1
% does not look nice, try deleting the line with the fontenc.


\title[Présentation de Projet RSE] % (optional, use only with long paper titles)
{Réseau social d'entreprise}

\subtitle
{ Formation développeur .NET pour des systèmes d'information d'entreprise} % (optional)

\author[Régis, Ekatherina] % (optional, use only with lots of authors)
{L.~Régis \and M.~Ekatherina}


\date[Short Occasion] % (optional)
{11/01/2019}

\AtBeginSubsection[]
{
  \begin{frame}<beamer>{Outline}
    \tableofcontents[currentsection,currentsubsection]
  \end{frame}
}


% If you wish to uncover everything in a step-wise fashion, uncomment
% the following command: 

%\beamerdefaultoverlayspecification{<+->}


\begin{document}

\begin{frame}
  \titlepage
\end{frame}

\begin{frame}{Table des matières}
  \tableofcontents
  % You might wish to add the option [pausesections]
\end{frame}


% Since this a solution template for a generic talk, very little can
% be said about how it should be structured. However, the talk length
% of between 15min and 45min and the theme suggest that you stick to
% the following rules:  

% - Exactly two or three sections (other than the summary).
% - At *most* three subsections per section.
% - Talk about 30s to 2min per frame. So there should be between about
%   15 and 30 frames, all told.

\section{Introduction}
\begin{frame}{L'équipe}
\begin{minipage}{0.45\textwidth}
Lebrun Régis
\begin{itemize}
  \item Baccalauréat 3ème année en Informatique industrielle (ESI)
  \item Formation .NET TechnofuturTIC
  \item Cours du soir Bachelier en informatique de gestion (EPHEC)
\end{itemize}
\end{minipage}%
\hfill
\begin{minipage}{0.45\textwidth}
\begin{tabular}{|p{\textwidth}}
Ekatherina Moskovkina
\begin{itemize}
\item D'origine Russe
\item Bachelier en chimie
\item Travaillé comme developpeur Python en Russie
\end{itemize}
\end{tabular}
\end{minipage}%
\end{frame}
\section{Objectifs}
\begin{frame}{Les objectifs du projet}
 Nous avons divisé notre projet en 8 parties:
  \begin{itemize}
  \item Employés
\item Départements
 \item Evènements d'entreprise
 \item Projets
 \item Equipes de projet
 \item Tâches d'équipe
  \item Messagerie instantanée
\item Partage de documents
    \end{itemize}
\end{frame}

\begin{frame}{Gestion des Employés}
\begin{figure}[h!]
  \includegraphics[width=0.2\textwidth]{images/employee}
\end{figure}
Objectifs:
  \begin{itemize}
    \item niveau d'acces
  \end{itemize}
\end{frame}

\begin{frame}{Gestion des Evènements}
\begin{figure}[h!]
  \includegraphics[width=0.20\textwidth]{images/Event}
\end{figure}
Objectifs:
  \begin{itemize}
    \item Inscription ouverte ou fermée
    \item Geolocalisation
  \end{itemize}
\end{frame}


\begin{frame}{Gestion des Projets}
\begin{figure}[h!]
  \includegraphics[width=0.20\textwidth]{images/manager_work_background_colleagues_meeting_icons_cartoon_characters_6836720}
\end{figure}
Objectifs:
  \begin{itemize}
    \item Gestion de temps (deadline)
    \item Divisé en équipes
    \item Chef de projet
  \end{itemize}
\end{frame}

\begin{frame}{Gestion des Equipes}
\begin{figure}[h!]
  \includegraphics[width=0.5\textwidth]{images/team}
\end{figure}
Objectifs:
  \begin{itemize}
    \item Chef d'équipe
    \item Gestion des tâches d'équipe
  \end{itemize}
\end{frame}

\begin{frame}{Gestion des Tâches}
\begin{figure}[h!]
  \includegraphics[width=0.3\textwidth]{images/tasks}
\end{figure}
Objectifs::
  \begin{itemize}
    \item Statut de tâche
    \item Tashes d'équipe et de projét
  \end{itemize}
\end{frame}

\begin{frame}{Gestion des Messages}
\begin{figure}[h!]
  \includegraphics[width=0.20\textwidth]{images/chat}
\end{figure}
Objectifs:
  \begin{itemize}
    \item Messages d'équipe, de projet, de tâche
    \item Messages personnels
  \end{itemize}
\end{frame}

\begin{frame}{Gestion des Documents}
\begin{figure}[h!]
  \includegraphics[width=0.15\textwidth]{images/document-open-hi}
\end{figure}
Objectifs:
  \begin{itemize}
    \item Partage de documents sur projets, événements, départements, équipes et tâches
  \end{itemize}
\end{frame}



\section{Structure}

\subsection{Structure du projet}
\def\titlename{Structure du projet}
\def\imagepath{images/dataflow.png}
\begin{frame}{\titlename}
  \begin{center}
    \includegraphics[width=0.9\textwidth,height=0.9\textheight,keepaspectratio]{\imagepath}
  \end{center}
\end{frame}

\subsection{Carte du site}
\def\titlename{Carte du site}
\def\imagepath{images/sitemap.png}
\begin{frame}{\titlename}
  \begin{center}
    \includegraphics[width=0.9\textwidth,height=0.9\textheight,keepaspectratio]{\imagepath}
  \end{center}
\end{frame}

\section{Ce que l'on a appris}
\begin{frame}{La base de données}

 La base de données:
  \begin{itemize}
  \item Vérification sur redondance d'informations permet de mitiger les bugs/modifications malveillantes
  \item Procédures stockées
  \item Soft delete
  \end{itemize}
\end{frame}
\begin{frame}{Mots de passe}
\begin{figure}[h!]
  \includegraphics[width=0.10\textwidth]{images/website_-_padlock-512}
\end{figure}
  Securité : 
  \begin{itemize}
  \item Stockage en clear - pas secure!
  \item SHA-256
  \item Technique du salage
  \end{itemize}
\end{frame}
\begin{frame}{Les Areas}
\begin{figure}[h!]
  \includegraphics[width=0.10\textwidth]{images/no__not__stop__wrong-512}
\end{figure}
  Les Areas permettent :
  \begin{itemize}
  \item La séparation totale entre différentes zones (Employe/Admin)
  \item Protection de controleurs avec des droits étendus afin d'empêcher toute faille de sécurité
  \end{itemize}
\end{frame}
\begin{frame}{Les Attributes}

  Les attributes peuvent servir à :\dots
  \begin{itemize}
  \item Filtrer les droits d'accès
  \item Définir la mise en forme des données (format dates)
\item Définir des contraintes d'intégrités
\item Définir des métadonnées en général (nom long des attributs)
  \end{itemize}
\end{frame}
\begin{frame}{Variables de session}

  You can create overlays\dots
  \begin{itemize}
  \item using the \texttt{pause} command:
    \begin{itemize}
    \item
      First item.
      \pause
    \item    
      Second item.
    \end{itemize}
  \end{itemize}
\end{frame}
\begin{frame}{Le stockage de documents}
\begin{figure}[h!]
  \includegraphics[width=0.10\textwidth]{images/iconfinder_website_-_folder_3440836}
\end{figure}
 Avantages du stockage sur base de données :
  \begin{itemize}
  \item Versioning facile
  \item Pas dans le système d'exploitation
  \item FILESTREAM
  \end{itemize}
\end{frame}
\begin{frame}{La messagerie instantanée}
\begin{figure}[h!]
  \includegraphics[width=0.10\textwidth]{images/-59-512}
\end{figure}

  SignalR permet de créer des messages instantanés
  \begin{itemize}
  \item Quand l'utilisateur crée un message, il notifie le hub
  \item Le hub notifie tout qui doit recevoir ce message
    \begin{itemize}
    \item Mais notre système d'authentification (par HttpContext) n'est pas accessible dans le hub
    \item Alors le hub notifie tous les utilisateurs qui sont sur la page avec le script, et les utilisateurs retrievent que les messages appropriés à partir de la base de données
    \end{itemize}
  \item JQuery permet d'ajouter nouveaux messages sans recharger toute la page
  \end{itemize}
\end{frame}

\section*{En résumé}

\begin{frame}{En résumé}

  % Keep the summary *very short*.
  \begin{itemize}
  \item
    The \alert{first main message} of your talk in one or two lines.
  \item
    The \alert{second main message} of your talk in one or two lines.
  \item
    Perhaps a \alert{third message}, but not more than that.
  \end{itemize}
  
  % The following outlook is optional.
  \vskip0pt plus.5fill
  \begin{itemize}
  \item
    Améliorations possibles
    \begin{itemize}
    \item
      Ajout de mappers pour les forms.
    \item
      Ajout d'un sel différent pour tous les utilisateurs.
    \end{itemize}
  \end{itemize}
\end{frame}

\section{Présentation vidéo}
\begin{frame}{Questions}
\begin{center}
\LARGE Présentation vidéo
 \end{center}
\end{frame}

\section{Questions}
\begin{frame}{Questions}
\begin{center}
\LARGE Avez-vous des questions?
 \end{center}
\end{frame}

\end{document}


